\subsection{Exercices}

\targetExo{fsec1_ex1}
\begin{exercicebox}[Exercice 1 : Profit d'une position Long Call]
Vous avez acheté un Call de Strike $100$ € pour une prime de $5$ €. À l'échéance, le sous-jacent vaut $112$ €. 
Calculez votre Payoff, puis déduisez-en votre profit net (en ignorant la valeur temporelle de la prime).
\navToCorr{fsec1_ex1}
\end{exercicebox}

\targetExo{fsec1_ex2}
\begin{exercicebox}[Exercice 2 : Profit d'une position Short Put]
Vous vendez un Put de Strike $K = 200$ € et encaissez immédiatement une prime de $10$ €. À l'échéance, l'action cote $185$ €. 
L'acheteur exerce-t-il son option contre vous ? Quel est votre profit (ou perte) net ?
\navToCorr{fsec1_ex2}
\end{exercicebox}

\targetExo{fsec1_ex3}
\begin{exercicebox}[Exercice 3 : Statut d'un Call (Moneyness)]
L'action Apple cote actuellement $175$ \$. Vous observez trois Calls européens sur le marché avec des Strikes différents : $K_1 = 160$ \$, $K_2 = 175$ \$, et $K_3 = 190$ \$. 
Déterminez le statut (ITM, OTM, ou ATM) de chacune de ces trois options aujourd'hui.
\navToCorr{fsec1_ex3}
\end{exercicebox}

\targetExo{fsec1_ex4}
\begin{exercicebox}[Exercice 4 : Combinaison de flux]
Vous construisez le portefeuille suivant aujourd'hui : vous achetez 1 action à $S_0 = 100$ € et vous achetez 1 Put de Strike $K = 100$ €.
Écrivez mathématiquement la valeur de ce portefeuille à l'échéance $T$ en fonction de $S_T$. Montrez que ce profil de flux ressemble fortement à un produit dérivé bien connu (Stratégie dite du Protective Put).
\navToCorr{fsec1_ex4}
\end{exercicebox}

\targetExo{fsec1_ex5}
\begin{exercicebox}[Exercice 5 : Stratégie Straddle (Achat de volatilité)]
Un investisseur anticipe un fort mouvement sur une action (actuellement à $100$ €), mais il ignore la direction. Il achète simultanément 1 Call de Strike $100$ (Prime = 5 €) et 1 Put de Strike $100$ (Prime = 5 €). Le coût total est donc de $10$ €.
Calculez le Payoff global de ce portefeuille si l'action finit à $120$ €. Calculez-le si l'action finit à $80$ €.
\navToCorr{fsec1_ex5}
\end{exercicebox}